\subsection{The process and its simulation}

The Hawkes process itself is \citet{hawkes1971}, which defines the self-exciting
and mutually exciting classes and derives the spectral density that gives the
process its name; the conditional-intensity and branching-matrix formulation used
throughout this paper is the modern reading of that construction. The thinning
simulation used here is \citet{ogata1981}, which is still the standard way to draw
exact samples for an inhomogeneous point process with a bounded intensity: one
simulates an upper bound and accepts each candidate with probability proportional
to the intensity ratio. The package implements that algorithm against a second,
independent branching simulator and compares the two at the distributional level,
so a disagreement between the simulators fails a test rather than a downstream
study. Neither of those works discusses the sampling distribution of a
maximum-likelihood fit, and neither is a claim this paper revises.

\subsection{Inference for Hawkes processes}

The methodological literature on inference is recent and largely asymptotic.
\citet{laub2015} is the standard reference for the likelihood, the exponential
kernel's recursions and the properties that make the estimator tractable, and it
sets out the heuristic for consistency in the stationary case. \citet{chen2013}
develops inference under a nonstationary specification, with a parametric
likelihood corrected for nonstationarity and an application to ultra-high-frequency
returns. Both establish consistency and asymptotic normality in their settings, and
both state the covariance in the sandwich form of \eqref{eq:sandwich}. Neither
measures what that covariance does at a fixed number of events, which is the
question this paper asks, and neither reports the ratio of mean reported standard
error to empirical standard deviation anywhere, because neither is a simulation
study of the estimator against itself.

The asymptotic results are not in dispute here. Section~\ref{sec:theory} uses them
as the starting point and asks what the information equality \eqref{eq:equality}
delivers in a finite window. The finding that the reported interval is too narrow at
$N = 2000$ and adequate at $N = 50\,000$ is consistent with those results rather
than in tension with them: an asymptotic approximation that is good at five
thousand events per type can be poor at two thousand, and the paper's contribution
is to locate where the crossing happens in a box it measured, not to argue that
asymptotics fail.

\subsection{Applications that report parameters without intervals}

The finance literature fits these models routinely and reports the parameters. The
vocabulary this paper uses, the branching matrix, its spectral radius and the
near-critical regime, is set by the review of \citet{bacry2015}, which is the
clearest statement of why $\rho(\Gamma)$ close to one is the regime of interest.
\citet{filimonov2012} estimate a Hawkes branching ratio on financial data and use
the ratio itself as a reflexivity diagnostic, with the fitted parameter as the
reported quantity and no sampling distribution attached to the maximum-likelihood
fit. \citet{bacry2013} model microstructure noise with mutually exciting processes
and report fitted parameters and goodness of fit rather than the coverage of their
intervals.

That body of work is not doing anything wrong, and this paper does not claim it is.
A diagnostic whose value is read as a level does not need an interval to be useful,
and a parameter reported to six figures can be a legitimate summary of a fitted
model. What those papers do not report is the property this paper measures: whether
the standard error that comes out of the same fit describes the spread of estimates
the procedure would produce. When the answer is that it understates that spread by
a third at the sample sizes those papers often work at, a reader who treats the
reported interval as one is reading a number with no measured coverage behind it.
The paper's claim is narrow and it is about the interval, not the parameter.

\subsection{Software and reproducibility}

The algorithmic treatment closest to the implementation here is \citet{bacry2015},
which also indexes the simulation and estimation tooling that practice relies on.
None of the software this paper could compare against ships a measured calibration
map for its own standard errors, and none refuses to report an interval where the
estimator is undefined; those are the two behaviours the package adds, and they are
additions to the inference layer rather than to the likelihood. The design of this
study, frozen before measurement and amended six times with each amendment dated
and justified, follows the practice the reproducibility literature argues for
without needing to cite it: the point is that a reader can determine what was known
when each decision was made, and the appendix is there so they can.

No speed or accuracy comparison with another library is made here. The capability
table the package ships times this implementation on one machine and states no
claim about any other.
